\section{Conclusion}
\label{sec:conclusion}

The Heisenberg pentamer chain connects a change of the local composite spin
to a quantitative redistribution of low-energy spectral weight.
In the isolated-cell ranges studied, positive and negative outer intra-cell
exchange select spin-$1/2$ and spin-$3/2$ ground multiplets, respectively.
Projecting the physical spins onto these multiplets fixes both the leading
antiferromagnetic inter-cell exchange and the momentum-dependent interference
factor. The latter favors different physical replicas of the same
effective-chain antiferromagnetic momentum: $q=\pi$ for $\alpha=0.5$ and
$q=0.6\pi$ for $\alpha=-0.5$.
Two-cell exact diagonalization verifies the resulting absolute low-band
weights and excitation scale in the specified weak-coupling regime.
Reflection parity independently explains why the lowest excited doublet of
the positive-$\alpha$ isolated cell is invisible to a $q=\pi$ spin probe.

Finite-chain calculations place these local predictions in a broader
coupling range. Positive-coupling static diagnostics support an effective
antiferromagnetic chain at the tested points, while negative couplings
produce a size-dependent reconstruction of the dominant static momentum.
The alternative trimer--dimer limit near $\alpha=0$ explains why the
isolated-quartet description need not persist there.
Complete connected $L=40$ spectra at $(\alpha,\beta)=(\pm0.5,0.5)$ compare
the responses at one common resolution, outside the controlled
weak-coupling regime. Together, the results distinguish a change in the
visibility of collective spin fluctuations from a demonstrated phase
transition or a unique assignment of high-energy bands.
